% Source problem: https://www.ap.t.u-tokyo.ac.jp/themes/ap_t_u-tokyo/graduate/pdf/h31ap_exam.pdf
\documentclass[11pt,a4paper]{article}
\usepackage[margin=24mm]{geometry}
\usepackage{fontspec}
\usepackage{xeCJK}
\setCJKmainfont{HaranoAjiMincho-Regular.otf}
\usepackage{amsmath,amssymb}
\setlength{\parindent}{1em}
\setlength{\parskip}{0.35em}
\title{物理工学専攻 院試解答集\\2019年度 物理学II 第2問}
\author{政岡凜太郎}
\date{}
\begin{document}
\maketitle

\section*{[1] 臨界点}
状態方程式を粒子数密度$n$で微分すると
\begin{align}
 \left(\frac{\partial P}{\partial n}\right)_T
 &=\frac{k_BT}{(1-bn)^2}-2an,\\
 \left(\frac{\partial^2P}{\partial n^2}\right)_T
 &=\frac{2bk_BT}{(1-bn)^3}-2a.
\end{align}
両方をゼロとし、式を割り合わせると
\begin{equation}
 \frac{2b}{1-bn_c}=\frac1{n_c},\qquad
 \boxed{n_c=\frac1{3b}}.
\end{equation}
第2式から臨界温度が得られ、状態方程式へ代入すれば臨界圧力は
\begin{equation}
 \boxed{T_c=\frac{8a}{27k_Bb}},\qquad
 \boxed{P_c=\frac{a}{27b^2}}.
\end{equation}

\section*{[2] 臨界密度での等温圧縮率}
固定した粒子数$N$について$n=N/V$なので、$(\partial V/\partial n)_N=-V/n$である。圧縮率の定義に連鎖律を使うと
\begin{equation}
 K_T=-\frac1V\left(\frac{\partial V}{\partial P}\right)_{T,N}
 =-\frac1V\left(-\frac Vn\right)
 \left(\frac{\partial P}{\partial n}\right)_T^{-1}
 =\left[n\left(\frac{\partial P}{\partial n}\right)_T\right]^{-1}.
\end{equation}
$n=n_c=1/(3b)$を代入すると
\begin{align}
 K_T&=\frac1{n_c}\left(\frac94k_BT-\frac{2a}{3b}\right)^{-1}\\
 &=\boxed{\frac{4}{9n_ck_B(T-T_c)}}\qquad (T>T_c),
\end{align}
ここで$2a/(3b)=(9/4)k_BT_c$を用いた。$T\downarrow T_c$で圧縮率が発散する。

\section*{[3] 理想気体の自由エネルギーと状態量}
運動量積分は
\begin{equation}
 \int_{-\infty}^{\infty}e^{-p^2/(2mk_BT)}\,dp=\sqrt{2\pi mk_BT}
\end{equation}
なので
\begin{equation}
 Z_0=\frac{V^N}{N!h^{3N}}(2\pi mk_BT)^{3N/2}.
\end{equation}
Stirling近似を使って$F_0=-k_BT\ln Z_0$を整理すると
\begin{equation}
 F_0=-Nk_BT\left[\ln\!\left(\frac{V}{N}\right)+1
 +\frac32\ln\!\left(\frac{2\pi mk_BT}{h^2}\right)\right].
\end{equation}
これを各変数で微分すると
\begin{align}
 P_0&=-\left(\frac{\partial F_0}{\partial V}\right)_{T,N}
 =\boxed{\frac{Nk_BT}{V}},\\
 S_0&=-\left(\frac{\partial F_0}{\partial T}\right)_{V,N}
 =\boxed{Nk_B\left[\ln\!\left(\frac{V}{N}\right)+\frac52
 +\frac32\ln\!\left(\frac{2\pi mk_BT}{h^2}\right)\right]},\\
 \mu_0&=\left(\frac{\partial F_0}{\partial N}\right)_{T,V}
 =\boxed{-k_BT\left[\ln\!\left(\frac{V}{N}\right)
 +\frac32\ln\!\left(\frac{2\pi mk_BT}{h^2}\right)\right]},\\
 U_0&=F_0+TS_0=\boxed{\frac32Nk_BT}.
\end{align}
対数は組合せをまとめると無次元であり、例えば
$\mu_0=k_BT\ln[n\lambda_T^3]$、$\lambda_T=h/\sqrt{2\pi mk_BT}$と書ける。

\section*{[4] 高温・低密度での自由エネルギー}
相互作用の積分は
\begin{equation}
 \int_\ell^\infty r^2\frac{u(r)}{k_BT}\,dr
 =-\frac{\epsilon\ell^3}{3k_BT}.
\end{equation}
また$4\pi\ell^3/3=2v$だから、与えられた近似式の角括弧は
$1+2\epsilon v/(k_BTV)$となる。$N$が大きく、密度が低く、
 $\epsilon/(k_BT)\ll1$として引力に由来する対数を展開し、$N-1\simeq N$とすれば
\begin{align}
 \ln A&\simeq N\ln\left(1-\frac{Nv}{V}\right)
 +\frac{N(N-1)}2\ln\left(1+\frac{2\epsilon v}{k_BTV}\right)\\
 &\simeq N\ln\left(1-\frac{Nv}{V}\right)
 +\frac{N^2\epsilon v}{k_BTV}.
\end{align}
排除体積因子は問題文の形のまま保ち、引力補正のみを一次で近似する。従って
\begin{equation}
 \boxed{F\simeq F_0-Nk_BT\ln\left(1-\frac{Nv}{V}\right)
 -\frac{N^2\epsilon v}{V}}.
\end{equation}
低密度でさらに展開すれば排除体積補正の先頭は$N^2vk_BT/V$となる。排除体積は自由エネルギーを増やし、引力は減らす。

\section*{[5] 圧力と [6] 定数の対応}
同じ近似で自由エネルギーを密度形に書くと
\begin{equation}
 F\simeq -Nk_BT\left[\ln\!\left(\frac{V-Nv}{Nh^3}\right)+1
 +\frac32\ln(2\pi mk_BT)+\frac{\epsilon vN}{k_BTV}\right].
\end{equation}
この形から体積微分を取ると
\begin{equation}
 P=-\left(\frac{\partial F}{\partial V}\right)_{T,N}
 =\frac{Nk_BT}{V-Nv}-\frac{\epsilon vN^2}{V^2}
 =\boxed{\frac{nk_BT}{1-vn}-\epsilon v n^2}.
\end{equation}
よって van der Waals 方程式と比較して
\begin{equation}
 \boxed{b=v,\qquad a=\epsilon v}.
\end{equation}
各項の次元は圧力であり、引力項の符号は負である。
\end{document}
